\documentclass[12pt]{article}
\usepackage{../t068report}
% Auto-generated by make_numbers.py from results/*.json — do not edit by hand
\newcommand{\nRowsRaw}{284,807}
\newcommand{\nFraudRaw}{492}
\newcommand{\nDup}{1,081}
\newcommand{\nDupFraud}{19}
\newcommand{\nRows}{283,726}
\newcommand{\nFraud}{473}
\newcommand{\fraudRate}{0.167\%}
\newcommand{\nSkewed}{13}
\newcommand{\nHeavy}{18}
\newcommand{\ksTop}{V14 (0.84), V10 (0.80), V12 (0.78), V4 (0.76), V11 (0.75), V17 (0.74)}
\newcommand{\trainRows}{170,236}
\newcommand{\trainFraud}{342}
\newcommand{\trainRate}{0.201\%}
\newcommand{\trainLegit}{169,894}
\newcommand{\valRows}{56,745}
\newcommand{\valFraud}{57}
\newcommand{\valRate}{0.100\%}
\newcommand{\valLegit}{56,688}
\newcommand{\testRows}{56,745}
\newcommand{\testFraud}{74}
\newcommand{\testRate}{0.130\%}
\newcommand{\testLegit}{56,671}
\newcommand{\selK}{4}
\newcommand{\selScaler}{quantile}
\newcommand{\selFS}{pca}
\newcommand{\selD}{28}
\newcommand{\perComp}{434}
\newcommand{\kBIC}{24}
\newcommand{\kAIC}{24}
\newcommand{\gmmTestAP}{0.686}
\newcommand{\gmmValAP}{0.695}
\newcommand{\gmmCIlo}{0.572}
\newcommand{\gmmCIhi}{0.808}
\newcommand{\gmmROC}{0.934}
\newcommand{\gmmRecallA}{75.7\%}
\newcommand{\gmmPrecA}{59.6\%}
\newcommand{\gmmAlertsA}{94}
\newcommand{\gmmTPA}{56}
\newcommand{\gmmFPA}{38}
\newcommand{\gmmFPRA}{0.067\%}
\newcommand{\gmmRecallB}{78.4\%}
\newcommand{\gmmPrecB}{26.2\%}
\newcommand{\gmmAlertsB}{221}
\newcommand{\gmmTPB}{58}
\newcommand{\gmmFPB}{163}
\newcommand{\gmmFPRB}{0.288\%}
\newcommand{\gmmRecallC}{78.4\%}
\newcommand{\gmmPrecC}{14.5\%}
\newcommand{\gmmAlertsC}{399}
\newcommand{\gmmTPC}{58}
\newcommand{\gmmFPC}{341}
\newcommand{\gmmFPRC}{0.602\%}
\newcommand{\gmmPA}{90\%}
\newcommand{\gmmPB}{56\%}
\newcommand{\gmmPC}{28\%}
\newcommand{\gmmPD}{12\%}
\newcommand{\feAblationRows}{%
pca & 0.296$\pm$0.290 & 0.578$\pm$0.168 & 0.085$\pm$0.010 & \textbf{0.698$\pm$0.009} \\
pca+amount & 0.246$\pm$0.133 & 0.165$\pm$0.121 & 0.453$\pm$0.114 & 0.659$\pm$0.021 \\
full & 0.507$\pm$0.291 & 0.468$\pm$0.212 & 0.096$\pm$0.024 & 0.631$\pm$0.086 \\
}
\newcommand{\selectionRows}{%
1 & 434 & 2.995 & -37.42 & 0.691$\pm$0.000 \\
2 & 869 & 2.218 & -27.71 & 0.691$\pm$0.001 \\
3 & 1,304 & 1.971 & -24.58 & 0.683$\pm$0.015 \\
4 & 1,739 & 1.795 & -22.38 & 0.698$\pm$0.012 \\
6 & 2,609 & 1.486 & -18.43 & 0.510$\pm$0.233 \\
8 & 3,479 & 1.251 & -15.42 & 0.596$\pm$0.107 \\
12 & 5,219 & 1.025 & -12.46 & 0.591$\pm$0.036 \\
16 & 6,959 & 0.819 & -9.60 & 0.591$\pm$0.153 \\
24 & 10,439 & 0.678 & -7.56 & 0.554$\pm$0.074 \\
}
\newcommand{\diagBestAP}{0.108}
\newcommand{\diagBestK}{24}
\newcommand{\mainRows}{%
GMM (K=4, full) & 0.686 & [0.57, 0.81] & 0.934 & 0.757 & 0.784 & 0.784 & 0.90 & 0.56 \\
GMM, BIC choice (K=24, full) & 0.626 & [0.51, 0.76] & 0.955 & 0.662 & 0.757 & 0.770 & 0.86 & 0.53 \\
Single Gaussian (K=1) & 0.691 & [0.58, 0.80] & 0.963 & 0.757 & 0.757 & 0.770 & 0.92 & 0.56 \\
Single Gaussian, standard scaling & 0.163 & [0.12, 0.22] & 0.959 & 0.081 & 0.514 & 0.716 & 0.18 & 0.28 \\
Isolation Forest & 0.050 & [0.04, 0.07] & 0.939 & 0.014 & 0.203 & 0.432 & 0.02 & 0.06 \\
Supervised LogReg (uses labels) & 0.739 & [0.65, 0.83] & 0.981 & 0.770 & 0.851 & 0.851 & 0.88 & 0.56 \\
}
\newcommand{\apKOne}{0.691}
\newcommand{\apStd}{0.163}
\newcommand{\apIso}{0.050}
\newcommand{\apSup}{0.739}
\newcommand{\apBIC}{0.626}
\newcommand{\stabilityRows}{%
4 & quantile (selected) & 0.705$\pm$0.032 & [0.65, 0.74] & 0.897 & 0.87 \\
4 & standard & 0.505$\pm$0.224 & [0.15, 0.74] & 0.851 & 0.77 \\
8 & quantile (selected) & 0.650$\pm$0.033 & [0.59, 0.69] & 0.866 & 0.79 \\
8 & standard & 0.418$\pm$0.226 & [0.11, 0.68] & 0.798 & 0.63 \\
16 & quantile (selected) & 0.604$\pm$0.089 & [0.49, 0.74] & 0.906 & 0.74 \\
16 & standard & 0.511$\pm$0.229 & [0.10, 0.72] & 0.813 & 0.63 \\
}
\newcommand{\stabQstd}{0.032}
\newcommand{\stabSstd}{0.224}
\newcommand{\stabSmin}{0.148}
\newcommand{\stabSmax}{0.743}
\newcommand{\stabQmin}{0.652}
\newcommand{\stabQmax}{0.737}
\newcommand{\opRows}{%
0.1\% & 94 & 12.3 & 56/74 & 0.596 & 0.067\% \\
0.5\% & 221 & 28.9 & 58/74 & 0.262 & 0.288\% \\
1.0\% & 399 & 52.1 & 58/74 & 0.145 & 0.602\% \\
}
\newcommand{\costBestAlerts}{61}
\newcommand{\costTotal}{2,763}
\newcommand{\costNone}{7,728}
\newcommand{\condMax}{2,467}


\begin{document}

\maketitlepage{Phase 1: Machine Learning Evaluation}{Gaussian Mixture density scoring of legitimate card transactions}

%==================================================================
\section*{Abstract}
\addcontentsline{toc}{section}{Abstract}
Card fraud is rare (\fraudRate{} of transactions in our data), its patterns change over time, and confirmed labels
arrive only weeks after a transaction. A practical fraud system is therefore a \emph{triage} tool. It ranks
transactions so that a small team of analysts reviews the most suspicious ones first. In this first phase we frame
fraud detection as \emph{novelty detection}. A Gaussian Mixture Model (GMM), fitted by the Expectation--Maximisation
(EM) algorithm, learns the probability density of \emph{legitimate} transactions only, without ever seeing a fraud
label, and the negative log-likelihood $-\log p(\mathbf{x})$ serves as the anomaly score.

On the ULB credit-card dataset (\nRowsRaw{} transactions) we use a strictly chronological, leakage-controlled protocol:
duplicates are removed before splitting, the train/validation/test blocks are ordered in time, models are fitted on
legitimate rows only, and thresholds are chosen on validation and frozen for the test period. We find three things:
\begin{enumerate}
\item \textbf{Feature engineering matters most.} A per-feature quantile (Gaussianising) transform is the decisive
design choice. It turns a single-Gaussian detector from PR-AUC \apStd{} into \apKOne{} and removes most of EM's
seed-to-seed instability.
\item \textbf{More components do not mean better detection.} AIC and BIC keep asking for more mixture components
($K=\kBIC$), yet $K=1$ to $4$ detect fraud as well or better, and large mixtures are unstable.
\item \textbf{The final model is a strong triage tool.} The selected GMM ($K=\selK$) reaches a test PR-AUC of
\gmmTestAP{} (95\% CI \gmmCIlo--\gmmCIhi). At a false-positive budget of 0.1\% of legitimate customers it catches
\gmmRecallA{} of test frauds, and \gmmPA{} of its top-50 alerts are fraud. This comes close to a supervised model that
uses labels (\apSup{}).
\end{enumerate}
We document the false-positive trade-offs, an illustrative cost analysis and the limitations that motivate the deep
learning phase.

\medskip
\textbf{Keywords:} fraud detection, anomaly detection, novelty detection, Gaussian
mixture model, expectation--maximisation, model selection, imbalanced evaluation, precision--recall.

\newpage
\tableofcontents
\newpage

%==================================================================
\section{Introduction}

\subsection{Background and motivation}
Every card payment must be approved within milliseconds, yet only a few hundred out of hundreds of thousands are
fraudulent. Banks cannot investigate every transaction. Instead, a fraud system raises \emph{alerts}, and a limited
team of analysts reviews them, blocks cards, or contacts customers~\cite{dalpozzolo2018}. The key question is
therefore not ``is this transaction fraud?'' but ``\emph{which few transactions should analysts look at first?}''.
This is a ranking (triage) problem.

Supervised classifiers need many labelled frauds. In practice labels are scarce, arrive late (after chargebacks),
and describe only \emph{past} fraud patterns~\cite{bolton2002}. Legitimate transactions, by contrast, are abundant and
available immediately. This motivates \emph{novelty detection}: learn a model of normal behaviour and flag departures
from it~\cite{chandola2009,pimentel2014}.

\subsection{Problem statement}
The project brief asks us to frame card fraud as anomaly detection. Models are fitted primarily on legitimate
transactions, and probability-density scores are compared with reconstruction error on a held-out mixed set. In
Phase~1 (machine learning) we must:
\begin{itemize}
\item fit a Gaussian Mixture Model on legitimate transactions, tune the component count with AIC and BIC, and use low
log-likelihood as the anomaly score;
\item use a chronological split, fit scaling and the GMM only on legitimate training rows, keep the fraud label hidden
during training, select thresholds on validation data, and report base rates;
\item report PR-AUC, recall at fixed false-positive budgets and precision at a review capacity, and \emph{not} rely on
accuracy;
\item identify the gap of the GMM that a deep-learning model (Phase~2) should address.
\end{itemize}

\subsection{Objectives}
\begin{enumerate}
\item Build a reproducible, leakage-free evaluation protocol shared by every model of both phases.
\item Understand the data: quality, imbalance, temporal structure, separability and how Gaussian it is.
\item Engineer features and preprocessing that suit a Gaussian density model.
\item Fit, select and validate the GMM, and compare it with meaningful reference baselines.
\item Translate scores into operating policies (thresholds, review workload, cost) for fraud analysts.
\item Measure the GMM's weaknesses quantitatively, to motivate Phase~2.
\end{enumerate}

\subsection{Contributions}
\begin{itemize}
\item A leakage-controlled protocol: deduplication, a chronological 60/20/20 split, legitimate-only fitting, and
thresholds frozen on validation.
\item A feature-engineering study showing that \emph{marginal Gaussianisation} (quantile transform) is what makes a
GMM work on heavy-tailed transaction features.
\item An analysis of AIC/BIC model selection at large sample size, where the criteria and detection quality disagree.
\item A quantified stability study of EM restarts, measuring the ``instability in high dimensions'' gap named in the
brief.
\item A triage-oriented evaluation with operating points, workload per hour and an illustrative cost analysis.
\end{itemize}

\subsection{Organisation of the report}
Section~\ref{sec:lit} reviews related work. Sections~\ref{sec:data} and~\ref{sec:fe} analyse the data and describe the
feature engineering. Section~\ref{sec:method} presents the GMM, the EM algorithm, model selection and the evaluation
protocol. Section~\ref{sec:results} reports the experiments. Section~\ref{sec:disc} discusses the findings and
limitations, and Section~\ref{sec:concl} concludes. The appendices cover reproducibility and team contributions.

%==================================================================
\section{Literature Review}\label{sec:lit}

\subsection{Fraud detection as a learning problem}
Bolton and Hand~\cite{bolton2002} review statistical fraud detection. They distinguish \emph{supervised} methods,
which learn from labelled past fraud, from \emph{unsupervised} methods, which flag departures from normal behaviour.
They argue that unsupervised methods are essential because fraudsters adapt and labels are delayed and incomplete.
Dal~Pozzolo~\emph{et~al.} released the dataset we use~\cite{dalpozzolo2015} and later described a realistic
fraud-detection system in which investigators review a fixed number of alerts per day~\cite{dalpozzolo2018}. In that
setting the relevant metric is precision among the top-$k$ alerts, which is why we report precision at a review
capacity. Carcillo~\emph{et~al.}~\cite{carcillo2021} showed that unsupervised outlier scores add information to
supervised fraud classifiers, which supports treating anomaly scores as triage signals.

\subsection{Anomaly and novelty detection}
Chandola~\emph{et~al.}~\cite{chandola2009} classify anomaly-detection techniques as statistical (density-based),
distance-based, clustering-based, classification-based and spectral. Statistical techniques assume that normal
instances occur in high-probability regions of a stochastic model, which is exactly the GMM score $-\log p(\mathbf{x})$.
Pimentel~\emph{et~al.}~\cite{pimentel2014} review \emph{novelty detection} (training on normal data only). They
identify GMM-based density estimation as the canonical probabilistic approach and note its sensitivity to
dimensionality and to the choice of model order. Isolation Forest~\cite{liu2008} isolates anomalies with random
axis-aligned splits and needs no density model; we use it as a non-density reference.

\subsection{Gaussian mixtures, EM and model order}
Dempster, Laird and Rubin~\cite{dempster1977} formalised the EM algorithm for maximum-likelihood estimation with
latent variables and proved that each iteration does not decrease the likelihood. McLachlan and Peel~\cite{mclachlan2000}
and Bishop~\cite{bishop2006} give the standard treatment of finite mixtures: covariance structures (full, diagonal,
tied), the singularity problem when a component collapses onto a few points, and covariance regularisation. To choose
the number of components, Akaike's information criterion (AIC)~\cite{akaike1974} and Schwarz's Bayesian information
criterion (BIC)~\cite{schwarz1978} penalise the log-likelihood by model size. BIC is consistent for the true order
\emph{when the model family is correct}, an assumption that heavy-tailed real data violates.

\subsection{Evaluation under extreme class imbalance}
Davis and Goadrich~\cite{davis2006} show that a ROC curve can look excellent while precision is poor. Saito and
Rehmsmeier~\cite{saito2015} show empirically that precision--recall (PR) plots are more informative than ROC plots on
imbalanced data. Our primary metric is therefore PR-AUC (average precision). We never report accuracy: a model that
flags nothing is 99.83\% ``accurate''.

\subsection{Summary and research gap}
\begin{table}[H]
\centering\small
\caption{Summary of the literature that shapes our Phase-1 design.}
\label{tab:lit}
\begin{tabularx}{\linewidth}{@{}l X X@{}}
\toprule
\textbf{Work} & \textbf{Contribution} & \textbf{How we use it}\\
\midrule
Bolton \& Hand~\cite{bolton2002} & Unsupervised fraud detection is needed: labels are scarce and late & Legitimate-only training protocol\\
Dal Pozzolo \emph{et al.}~\cite{dalpozzolo2015,dalpozzolo2018} & Dataset; analysts review a fixed number of alerts & Precision at review capacity (P@$k$)\\
Dempster \emph{et al.}~\cite{dempster1977} & EM algorithm for latent-variable models & GMM fitting; local optima studied\\
Akaike~\cite{akaike1974}, Schwarz~\cite{schwarz1978} & AIC / BIC model selection & Sweep of $K$ and covariance type\\
Pimentel \emph{et al.}~\cite{pimentel2014} & GMM novelty detection is sensitive to dimension and order & The ``gap'' we measure\\
Liu \emph{et al.}~\cite{liu2008} & Isolation Forest & Non-density reference baseline\\
Saito \& Rehmsmeier~\cite{saito2015} & PR beats ROC under imbalance & PR-AUC as the primary metric\\
\bottomrule
\end{tabularx}
\end{table}
Most published work on this dataset trains \emph{supervised} classifiers with random (not time-ordered) splits and
reports accuracy or ROC-AUC. Fewer studies evaluate a label-free density model under a chronological, leakage-free
protocol with triage metrics, or examine how preprocessing and model order affect the stability of EM. That is the
focus of this phase.

%==================================================================
\section{Dataset and Exploratory Data Analysis}\label{sec:data}

\subsection{Dataset description}
We use the Credit Card Fraud Detection dataset released by the ULB Machine Learning Group~\cite{dalpozzolo2015}
(distributed via Kaggle). It contains \nRowsRaw{} transactions made by European cardholders over two days in
September~2013, of which \nFraudRaw{} are fraudulent. Table~\ref{tab:features} describes the columns. The notebooks
download the file from a public mirror and verify its SHA-256 checksum, so every run uses identical data.

\begin{table}[H]
\centering\small
\caption{Dataset columns.}
\label{tab:features}
\begin{tabularx}{\linewidth}{@{}l l X@{}}
\toprule
\textbf{Column} & \textbf{Type} & \textbf{Meaning}\\
\midrule
\texttt{Time} & numeric & Seconds elapsed since the first transaction in the file (0 to 172{,}792 s $\approx$ 48 h)\\
\texttt{V1}--\texttt{V28} & numeric & Principal components of confidential original features (anonymised by the bank)\\
\texttt{Amount} & numeric & Transaction amount in euros\\
\texttt{Class} & binary & 1 = fraud, 0 = legitimate (used only for evaluation and to remove frauds from training)\\
\bottomrule
\end{tabularx}
\end{table}

\subsection{Data quality}
The data has no missing values and all columns are numeric, so no imputation is needed. We found \nDup{} exact
duplicate rows, \nDupFraud{} of them frauds. We drop them \emph{before} splitting, because an identical row in both the
training and the test block leaks information and inflates test scores. After deduplication \nRows{} rows remain,
including \nFraud{} frauds.

\subsection{Class imbalance}
Fraud makes up only \fraudRate{} of the rows, about one transaction in 600 (Fig.~\ref{fig:balance}). A model that never
flags anything would be 99.83\% accurate and would catch no fraud. Accuracy is therefore meaningless here. We use the
triage metrics defined in Section~\ref{sec:metrics}.

\begin{figure}[H]
\centering
\includegraphics[width=0.48\linewidth]{01_class_balance.png}
\caption{Class balance on a logarithmic scale: legitimate transactions outnumber frauds by about 600 to 1.}
\label{fig:balance}
\end{figure}

\subsection{Time of day and transaction amount}
Transaction volume collapses in the early hours of the day, while the \emph{fraud rate} peaks there, at up to about
nine times the average (Fig.~\ref{fig:time}). The raw \texttt{Time} value only encodes a row's position within the
two-day file, so we keep the hour-of-day signal instead (Section~\ref{sec:fe}). \texttt{Amount} is extremely
right-skewed (skewness $>15$). Frauds have a lower median amount than legitimate transactions, because many are
small ``test'' charges, but they also have a long tail (Fig.~\ref{fig:amount}). Amount alone cannot separate the
classes.

\begin{figure}[H]
\centering
\includegraphics[width=0.82\linewidth]{01_time_of_day.png}
\caption{Transaction volume (bars) and fraud rate (line) by hour of day. Fraud is rare in absolute terms but
proportionally much more frequent at night.}
\label{fig:time}
\end{figure}

\begin{figure}[H]
\centering
\includegraphics[width=0.95\linewidth]{01_amount.png}
\caption{Distribution of the transaction amount (left, clipped at 1{,}000) and of $\log(1+\text{Amount})$ (right) for
legitimate and fraudulent transactions.}
\label{fig:amount}
\end{figure}

\subsection{Which components separate the classes?}
Labels are used here \emph{only for analysis}; the models never see them during training. We rank features by the
two-sample Kolmogorov--Smirnov (KS) statistic between the legitimate and the fraud distributions
(Fig.~\ref{fig:ks}). The strongest separators are \ksTop{}. Fig.~\ref{fig:top} shows that frauds sit in the
\emph{tails} of these components, far from where legitimate transactions concentrate. This is exactly the structure a
density-based detector exploits: a transaction in a low-density region of ``normal'' space is suspicious.

\begin{figure}[H]
\centering
\includegraphics[width=\linewidth]{01_ks_separation.png}
\caption{Kolmogorov--Smirnov separation between legitimate and fraud distributions, per feature (higher = more
different).}
\label{fig:ks}
\end{figure}

\begin{figure}[H]
\centering
\includegraphics[width=\linewidth]{01_top_features.png}
\caption{Densities of the six most separating components. Frauds (red) occupy the tails of the legitimate
distribution (blue).}
\label{fig:top}
\end{figure}

\subsection{Is legitimate behaviour Gaussian?}
A Gaussian density assumes symmetric, light-tailed marginals (skewness 0, excess kurtosis 0). On legitimate rows,
\nSkewed{} of the 31 model features have $|\text{skewness}|>1$ and \nHeavy{} have excess kurtosis above 3. The extreme
case is V28, with excess kurtosis of about 966. The Q--Q plots in Fig.~\ref{fig:qq} bend strongly away from the
Gaussian line. A single Gaussian is therefore a poor fit, and even mixture components are strained. This quantifies
the gap named in the brief: the Gaussian-mixture assumption may not hold for the legitimate transaction manifold.

\begin{figure}[H]
\centering
\includegraphics[width=\linewidth]{01_qq_plots.png}
\caption{Q--Q plots of legitimate transactions against a normal distribution. Points far from the line indicate heavy
tails and skew.}
\label{fig:qq}
\end{figure}

\subsection{Correlation structure}
Because V1--V28 are principal components, they are nearly uncorrelated over the whole dataset. \emph{Within the fraud
class}, however, they are strongly correlated (Fig.~\ref{fig:corr}): fraud moves several components together. A
full-covariance density model can recognise such an unusual \emph{combination} even when no single feature is
extreme, which a per-feature threshold rule could not.

\begin{figure}[H]
\centering
\includegraphics[width=0.95\linewidth]{01_correlation.png}
\caption{Correlation of V1--V18 among legitimate transactions (left) and among frauds (right).}
\label{fig:corr}
\end{figure}

\subsection{Chronological split and base rates}
Rows are ordered by \texttt{Time} and cut into three consecutive blocks (Table~\ref{tab:split}, Fig.~\ref{fig:split}).
Models are trained on the past and evaluated on the future, as they would be in production. The scaler and the GMM
are fitted only on the \trainLegit{} legitimate training rows. The fraud label is used once, to remove frauds from that
block. The validation block (mixed) is used to choose hyper-parameters and thresholds. The test block (mixed) is used
once, at the end. Base rates drift between blocks, so every precision value must be read against its block's base
rate.

\begin{table}[H]
\centering\small
\caption{Chronological split and base rates.}
\label{tab:split}
\begin{tabularx}{\linewidth}{@{}l r r r X@{}}
\toprule
\textbf{Block} & \textbf{Rows} & \textbf{Frauds} & \textbf{Fraud rate} & \textbf{Used for}\\
\midrule
Train (first 60\%) & \trainRows & \trainFraud & \trainRate & fitting scaler and GMM (legitimate rows only)\\
Validation (next 20\%) & \valRows & \valFraud & \valRate & choosing hyper-parameters and thresholds\\
Test (last 20\%) & \testRows & \testFraud & \testRate & final evaluation, used once\\
\bottomrule
\end{tabularx}
\end{table}

\begin{figure}[H]
\centering
\includegraphics[width=0.95\linewidth]{01_split_timeline.png}
\caption{Transactions per 30 minutes, coloured by block; red ticks mark frauds.}
\label{fig:split}
\end{figure}

\textbf{Limitation.} The validation and test blocks (hours 33--48) cover the daytime and evening of day two. The night-time hours, where the
fraud \emph{rate} peaks, occur only in the training block, so the time-of-day feature is evaluated on daytime
transactions only.

%==================================================================
\section{Feature Engineering}\label{sec:fe}

\subsection{Engineered features}
Table~\ref{tab:fe-design} lists every transformation and its rationale. The hour of day is placed on the unit circle,
so that 23:00 and 00:00 are neighbours:
\begin{equation}
h=\frac{\texttt{Time}\bmod 86400}{3600},\qquad \text{hour\_sin}=\sin\!\left(\frac{2\pi h}{24}\right),\qquad
\text{hour\_cos}=\cos\!\left(\frac{2\pi h}{24}\right).
\end{equation}

\begin{table}[H]
\centering\small
\caption{Feature-engineering decisions.}
\label{tab:fe-design}
\begin{tabularx}{\linewidth}{@{}l l X@{}}
\toprule
\textbf{Item} & \textbf{Choice} & \textbf{Reason}\\
\midrule
Duplicates & dropped before splitting & avoid train/test leakage\\
\texttt{Time} & $\rightarrow$ hour\_sin, hour\_cos; raw value dropped & keep the night-time fraud signal, remove the position-in-file artefact\\
\texttt{Amount} & $\rightarrow \log(1+\text{Amount})$ & remove extreme skew so one feature cannot dominate\\
V1--V28 & kept & already decorrelated principal components\\
Scaling & fitted on legitimate training rows only & no information from validation, test or fraud\\
Labels & used only to drop training frauds, set thresholds, score test & ``keep the fraud label hidden during unsupervised training''\\
\bottomrule
\end{tabularx}
\end{table}

We compare three feature sets: \emph{pca} (V1--V28, $d{=}28$), \emph{pca+amount} ($d{=}29$) and \emph{full}
($d{=}31$, adding the hour features).

\subsection{Scaling options}
Four scalers are compared, each fitted on legitimate training rows only:
\begin{itemize}
\item \textbf{Standard:} $z=(x-\mu)/\sigma$.
\item \textbf{Robust:} $z=(x-\text{median})/\text{IQR}$, which is less sensitive to outliers.
\item \textbf{Power (Yeo--Johnson):} a parametric transform that reduces skew, followed by standardisation.
\item \textbf{Quantile:} $z=\Phi^{-1}\big(\hat F(x)\big)$, where $\hat F$ is the empirical CDF of the feature on
legitimate training rows and $\Phi$ is the standard normal CDF. Every \emph{marginal} becomes exactly standard normal
on the training data. Values beyond the training range are clipped to the extreme quantiles.
\end{itemize}

%==================================================================
\section{Methodology}\label{sec:method}

\subsection{Overview}
Fig.~\ref{fig:pipeline} summarises the pipeline. Everything to the left of the scoring step uses legitimate training
rows only.

\begin{figure}[H]
\centering
\begin{tikzpicture}[node distance=0.55cm and 0.5cm,
  box/.style={draw,align=center,font=\small,minimum height=1.25cm,text width=2.25cm},
  hot/.style={box},
  arr/.style={-{Stealth[length=2mm]}}]
\node[box] (data) {Raw data\\\footnotesize 284{,}807 rows};
\node[box,right=of data] (split) {Deduplicate +\\time split\\\footnotesize 60/20/20};
\node[box,right=of split] (fe) {Features +\\quantile scaling\\\footnotesize legit train only};
\node[box,right=of fe] (gmm) {GMM by EM\\\footnotesize $K$ by AIC/BIC\\+ validation};
\node[hot,right=of gmm] (score) {Score\\$-\log p(\mathbf{x})$};
\node[hot,below=0.9cm of score] (queue) {Review queue\\\footnotesize threshold\\from validation};
\draw[arr] (data) -- (split); \draw[arr] (split) -- (fe); \draw[arr] (fe) -- (gmm); \draw[arr] (gmm) -- (score);
\draw[arr] (score) -- (queue);
\end{tikzpicture}
\caption{Phase-1 pipeline. The model sees only legitimate training transactions; labels are used only for choosing
thresholds (validation) and for the final evaluation (test).}
\label{fig:pipeline}
\end{figure}

\subsection{Gaussian Mixture Model}
A GMM with $K$ components models the density of legitimate transactions as
\begin{equation}
p(\mathbf{x})=\sum_{k=1}^{K}\pi_k\,\mathcal{N}(\mathbf{x}\mid\bm\mu_k,\bm\Sigma_k),\qquad \pi_k\ge 0,\quad \sum_{k=1}^{K}\pi_k=1,
\label{eq:gmm}
\end{equation}
\begin{equation}
\mathcal{N}(\mathbf{x}\mid\bm\mu,\bm\Sigma)=(2\pi)^{-d/2}\,|\bm\Sigma|^{-1/2}\exp\!\Big(-\tfrac12(\mathbf{x}-\bm\mu)^{\top}\bm\Sigma^{-1}(\mathbf{x}-\bm\mu)\Big).
\end{equation}
Each component can be read as one ``kind'' of normal transaction. The mixture can represent multi-modal and skewed
shapes, but every component is an ellipsoid.

\subsection{Expectation--Maximisation}
Given $N$ legitimate training rows, the log-likelihood
$\ell(\theta)=\sum_{i=1}^{N}\log\sum_{k=1}^{K}\pi_k\,\mathcal{N}(\mathbf{x}_i\mid\bm\mu_k,\bm\Sigma_k)$ has no closed-form
maximiser, because the logarithm of a sum couples the components. EM introduces the unobserved component membership
of each row and alternates two steps (Algorithm~\ref{alg:em}).

\textbf{E-step.} Compute the \emph{responsibility} of component $k$ for row $i$:
\begin{equation}
\gamma_{ik}=\frac{\pi_k\,\mathcal{N}(\mathbf{x}_i\mid\bm\mu_k,\bm\Sigma_k)}{\sum_{j=1}^{K}\pi_j\,\mathcal{N}(\mathbf{x}_i\mid\bm\mu_j,\bm\Sigma_j)}.
\end{equation}

\textbf{M-step.} With $N_k=\sum_i\gamma_{ik}$, re-estimate the parameters as responsibility-weighted averages:
\begin{equation}
\pi_k=\frac{N_k}{N},\qquad \bm\mu_k=\frac{1}{N_k}\sum_{i}\gamma_{ik}\mathbf{x}_i,\qquad
\bm\Sigma_k=\frac{1}{N_k}\sum_{i}\gamma_{ik}(\mathbf{x}_i-\bm\mu_k)(\mathbf{x}_i-\bm\mu_k)^{\top}+\lambda\mathbf{I}.
\end{equation}
Each EM iteration never decreases $\ell(\theta)$~\cite{dempster1977}, but EM converges only to a \emph{local} optimum
that depends on the initialisation. We initialise with k-means++ and study the effect of the random seed
(Section~\ref{sec:stab}). The ridge term $\lambda\mathbf{I}$ (\texttt{reg\_covar}, default $10^{-4}$) keeps every
covariance invertible when a component collapses onto a few nearly identical points; without it the likelihood can
diverge to infinity.

\begin{algorithm}[H]
\caption{EM for a Gaussian mixture fitted on legitimate transactions}
\label{alg:em}
\begin{algorithmic}[1]
\Require legitimate training rows $\mathbf{x}_1,\dots,\mathbf{x}_N$; number of components $K$; regulariser $\lambda$
\State Initialise $\bm\mu_k$ with k-means++; set $\pi_k=1/K$ and $\bm\Sigma_k$ from the k-means clusters
\Repeat
  \State \textbf{E-step:} $\gamma_{ik}\gets \pi_k\mathcal{N}(\mathbf{x}_i\mid\bm\mu_k,\bm\Sigma_k)\,/\,\sum_j\pi_j\mathcal{N}(\mathbf{x}_i\mid\bm\mu_j,\bm\Sigma_j)$
  \State \textbf{M-step:} $N_k\gets\sum_i\gamma_{ik}$;\; $\pi_k\gets N_k/N$;\; $\bm\mu_k\gets\frac{1}{N_k}\sum_i\gamma_{ik}\mathbf{x}_i$
  \State \hphantom{\textbf{M-step:}} $\bm\Sigma_k\gets\frac{1}{N_k}\sum_i\gamma_{ik}(\mathbf{x}_i-\bm\mu_k)(\mathbf{x}_i-\bm\mu_k)^{\top}+\lambda\mathbf{I}$
\Until{the increase of $\ell(\theta)$ is below a tolerance (or 500 iterations)}
\State \Return $\{\pi_k,\bm\mu_k,\bm\Sigma_k\}$; anomaly score $s(\mathbf{x})=-\log p(\mathbf{x})$
\end{algorithmic}
\end{algorithm}

\subsection{Anomaly score and decision rule}
The anomaly score is the negative log-likelihood under the legitimate density,
\begin{equation}
s(\mathbf{x})=-\log p(\mathbf{x}),
\end{equation}
so a transaction in a low-density region of legitimate behaviour receives a high score. A deployed system raises an
alert when $s(\mathbf{x})>t$. The threshold $t$ is the $(1-f)$-quantile of the scores of legitimate \emph{validation}
transactions, where $f$ is the false-positive budget. It is then frozen and applied to the test period. For $K=1$ the
score reduces (up to constants) to the squared Mahalanobis distance
$(\mathbf{x}-\bm\mu)^{\top}\bm\Sigma^{-1}(\mathbf{x}-\bm\mu)$, which we use as a baseline.

\subsection{Model selection with AIC and BIC}
With $p$ free parameters, maximised likelihood $\hat L$ and $n$ rows,
\begin{equation}
\mathrm{AIC}=2p-2\log\hat L,\qquad \mathrm{BIC}=p\log n-2\log\hat L .
\end{equation}
Both reward fit and penalise complexity. BIC penalises more strongly ($\log n>2$ for $n>7$), so it prefers smaller
models. Table~\ref{tab:params} gives the parameter counts.

\begin{table}[H]
\centering\small
\caption{Free parameters of a $K$-component GMM in $d$ dimensions.}
\label{tab:params}
\begin{tabular}{@{}l l l@{}}
\toprule
\textbf{Covariance} & \textbf{Free parameters $p$} & \textbf{Per component, $d=\selD$}\\
\midrule
Full & $(K-1)+Kd+K\,d(d+1)/2$ & \perComp{}\\
Diagonal & $(K-1)+2Kd$ & $2d=\the\numexpr 2*\selD\relax$\\
\bottomrule
\end{tabular}
\end{table}

\subsection{Computational complexity}
One EM iteration with full covariances costs $\mathcal{O}(N K d^2)$ for the E-step, which evaluates every Mahalanobis
form, plus $\mathcal{O}(K d^3)$ for the Cholesky factorisations. For $N\approx1.7\times10^{5}$, $d=\selD$ and $K\le24$
a fit takes seconds on a laptop CPU, but more than a minute on Colab's free CPU. The selection studies (about 100
fits) therefore use a random 40{,}000-row \emph{development subsample} of the legitimate training rows. The final model
and all baselines use all rows.

\subsection{Reference baselines}
To put the GMM's numbers in context we evaluate:
\begin{itemize}
\item a \textbf{single Gaussian} ($K=1$, Mahalanobis distance) on the same inputs;
\item the same single Gaussian on \textbf{standard-scaled} inputs, which isolates the effect of feature engineering;
\item the \textbf{BIC-selected} mixture;
\item \textbf{Isolation Forest}~\cite{liu2008} (300 trees), fitted on legitimate rows only;
\item a \textbf{supervised logistic regression} (class-balanced, $C=0.1$) that \emph{uses fraud labels}. This is an
upper reference showing what labels are worth, not a fair competitor.
\end{itemize}

\subsection{Evaluation protocol and metrics}\label{sec:metrics}
All scores follow the convention ``higher = more suspicious''. Let $y_i\in\{0,1\}$ be the label and $s_i$ the score.
\begin{itemize}
\item \textbf{PR-AUC (average precision).} $\text{AP}=\sum_n (R_n-R_{n-1})P_n$ over all thresholds, where $P_n$ and
$R_n$ are precision and recall. A random ranking scores the base rate (\testRate{} on test). We add a 95\% stratified
bootstrap confidence interval (1{,}000 resamples).
\item \textbf{Recall at a false-positive budget $f\in\{0.1\%,0.5\%,1\%\}$.} The threshold is set on validation so
that a fraction $f$ of legitimate validation transactions would be flagged. We report test recall and the
\emph{realised} test false-positive rate.
\item \textbf{Precision at a review capacity $k\in\{50,100,200,500\}$ (P@$k$).} The share of fraud among the $k$
highest-scored test transactions.
\item \textbf{ROC-AUC} as a secondary metric only.
\end{itemize}
Because EM depends on its initialisation, every comparison during development is the mean over three seeds.

\subsection{Experimental setup}
\begin{table}[H]
\centering\small
\caption{Experimental setup and hyper-parameters.}
\label{tab:setup}
\begin{tabularx}{\linewidth}{@{}l X@{}}
\toprule
\textbf{Item} & \textbf{Setting}\\
\midrule
Feature-set / scaler grid & \{pca, pca+amount, full\} $\times$ \{standard, robust, power, quantile\}; GMM full, $K=4$, 3 seeds\\
$K$ grid & \{1, 2, 3, 4, 6, 8, 12, 16, 24\}; full (3 seeds) and diagonal (1 seed) covariances\\
EM & k-means++ initialisation, $\lambda=10^{-4}$, max.\ 500 iterations; final model with 5 restarts (\texttt{n\_init=5})\\
Selection rule & AIC/BIC reported; $K$ chosen by mean validation PR-AUC\\
Stability study & 5 EM seeds for $K\in\{4,8,16\}$, quantile vs.\ standard inputs\\
Bootstrap & 1{,}000 stratified resamples for confidence intervals\\
Software & Python 3, scikit-learn, NumPy, pandas, SciPy, matplotlib~\cite{pedregosa2011}\\
Hardware & Google Colab (standard CPU runtime); global seed 42\\
\bottomrule
\end{tabularx}
\end{table}

%==================================================================
\section{Experiments and Results}\label{sec:results}

\subsection{Feature-engineering ablation}
Table~\ref{tab:feabl} and Fig.~\ref{fig:feabl} compare feature sets and scalers with a fixed GMM (full covariance,
$K=4$). \textbf{Quantile scaling is the best scaler for every feature set, and it has by far the smallest seed-to-seed
spread.} With standard, robust or power scaling, the heavy tails of legitimate features (Section~\ref{sec:data}) let a
few extreme legitimate values dominate the covariance estimates. Different EM runs then lock onto different extreme
points, and validation PR-AUC swings by $\pm0.1$--$0.3$. The selected inputs are \emph{\selFS{} + \selScaler{}}. With
only 57 validation frauds, the differences between feature sets under quantile scaling are within the noise level;
the decision that matters is the scaler.

\begin{table}[H]
\centering\small
\caption{Feature ablation: validation PR-AUC (mean $\pm$ std over 3 EM seeds; GMM, full covariance, $K=4$). Best in bold.}
\label{tab:feabl}
\begin{tabular}{@{}l c c c c@{}}
\toprule
\textbf{Feature set} & \textbf{standard} & \textbf{robust} & \textbf{power} & \textbf{quantile}\\
\midrule
\feAblationRows
\bottomrule
\end{tabular}
\end{table}

\begin{figure}[H]
\centering
\includegraphics[width=0.7\linewidth]{02_feature_ablation.png}
\caption{Feature ablation heat map (validation PR-AUC, mean $\pm$ std over 3 seeds).}
\label{fig:feabl}
\end{figure}

\subsection{Choosing $K$: information criteria versus detection quality}
Fig.~\ref{fig:sel} and Table~\ref{tab:sel} show the sweep on the selected inputs.

AIC, BIC and the label-free held-out log-likelihood \emph{all keep improving} up to the largest model tried
($K=\kBIC$). With tens of thousands of rows, the penalty $p\log n$ is small next to the likelihood gained by
describing the heavy tails ever more finely. Detection quality does not follow. Validation PR-AUC is flat or lower for
large $K$, and its seed-to-seed spread grows, because more components mean more local optima and extra components
begin to model rare legitimate behaviour in the sparse regions where fraud also lives. Diagonal covariances are far
worse at every $K$ (best \diagBestAP{} at $K=\diagBestK$): fraud is visible in \emph{combinations} of components,
which a diagonal model cannot represent.

We therefore report the BIC choice as a baseline, but select $K$ by mean validation PR-AUC, which is the purpose of the
validation block. The final model (full covariance, $K=\selK$) is refit on \emph{all} legitimate training rows with five
EM restarts, keeping the restart with the highest likelihood (a label-free choice).

\begin{figure}[H]
\centering
\includegraphics[width=\linewidth]{02_gmm_model_selection.png}
\caption{Model selection on quantile-scaled inputs. Left: AIC and BIC keep decreasing with $K$. Centre: held-out
legitimate log-likelihood keeps increasing. Right: validation PR-AUC (mean $\pm$ std over 3 seeds) does not improve
with $K$.}
\label{fig:sel}
\end{figure}

\begin{table}[H]
\centering\small
\caption{Full-covariance sweep on the development subsample (BIC in millions; PR-AUC mean $\pm$ std over 3 seeds).}
\label{tab:sel}
\begin{tabular}{@{}r r r r c@{}}
\toprule
\textbf{$K$} & \textbf{Parameters} & \textbf{BIC} & \textbf{Held-out $\ell$/row} & \textbf{Validation PR-AUC}\\
\midrule
\selectionRows
\bottomrule
\end{tabular}
\end{table}

\subsection{Main results}
Table~\ref{tab:main} reports test performance with thresholds fixed on validation, and Fig.~\ref{fig:pr} shows the
curves.

\begin{table}[H]
\centering\footnotesize
\caption{Test results (thresholds for the recall columns chosen on validation and frozen). P@$k$ = precision of the
top-$k$ alerts. Base rate on test = \testRate.}
\label{tab:main}
\setlength{\tabcolsep}{3pt}
\resizebox{\linewidth}{!}{%
\begin{tabular}{@{}l c c c c c c c c@{}}
\toprule
\textbf{Model} & \textbf{PR-AUC} & \textbf{95\% CI} & \textbf{ROC-AUC} & \textbf{Rec@0.1\%} & \textbf{Rec@0.5\%} & \textbf{Rec@1\%} & \textbf{P@50} & \textbf{P@100}\\
\midrule
\mainRows
\bottomrule
\end{tabular}}
\end{table}

The main observations are:
\begin{itemize}
\item The selected GMM reaches \textbf{PR-AUC \gmmTestAP{}}. At the strictest budget (0.1\% of legitimate customers
flagged on validation) it catches \textbf{\gmmRecallA{}} of test frauds with \gmmAlertsA{} alerts, of which \gmmPrecA{}
are fraud. Its top-50 queue is \textbf{\gmmPA{} fraud}.
\item \textbf{Mixture size barely matters:} a single Gaussian on the same inputs scores \apKOne{}; the BIC-selected
mixture ($K=\kBIC$) scores \apBIC{}.
\item \textbf{The scaler matters enormously:} the same single Gaussian on standard-scaled inputs collapses to
\apStd{}.
\item Isolation Forest reaches only \apIso{}. The label-using logistic regression (upper reference) reaches \apSup{}, so
the label-free GMM recovers most of what labels are worth.
\end{itemize}

ROC-AUC is high (0.93--0.98) for almost every model, including those with poor precision. This confirms that ROC is
uninformative under this imbalance~\cite{saito2015}.

\begin{figure}[H]
\centering
\includegraphics[width=\linewidth]{02_pr_and_budget_curves.png}
\caption{Test precision--recall curves (left) and recall versus false-positive rate on a log scale (right). Dotted
lines mark the 0.1\%, 0.5\% and 1\% budgets.}
\label{fig:pr}
\end{figure}

\subsection{Score distribution}
Fig.~\ref{fig:scores} shows the distribution of the GMM score on the test set. Most frauds receive scores far beyond
the bulk of legitimate transactions, while a minority overlap with legitimate behaviour. These overlapping frauds
cannot be caught without flagging many legitimate customers.

\begin{figure}[H]
\centering
\includegraphics[width=0.8\linewidth]{02_score_distribution.png}
\caption{GMM anomaly score on test (log density axis) with the three validation-chosen thresholds.}
\label{fig:scores}
\end{figure}

\subsection{Threshold analysis and false-positive trade-offs}
Table~\ref{tab:op} turns the scores into operating policies for the 7.7-hour test window. Tightening the budget from
1\% to 0.1\% shrinks the queue from \gmmAlertsC{} to \gmmAlertsA{} alerts while losing only a couple of frauds. Most
frauds are either very anomalous (caught at any budget) or look entirely normal (missed at every budget), so the
recall curve is flat between budgets (Fig.~\ref{fig:tradeoff}). The realised test false-positive rate stays below its
validation target because the test period is slightly ``cleaner'' than the validation period.

\begin{table}[H]
\centering\small
\caption{GMM operating points on the test period (thresholds from validation).}
\label{tab:op}
\begin{tabular}{@{}l r r r r r@{}}
\toprule
\textbf{FP budget} & \textbf{Alerts} & \textbf{Alerts / hour} & \textbf{Frauds caught} & \textbf{Precision} & \textbf{Realised FPR}\\
\midrule
\opRows
\bottomrule
\end{tabular}
\end{table}

\begin{figure}[H]
\centering
\includegraphics[width=0.8\linewidth]{02_threshold_tradeoff.png}
\caption{Precision and recall versus review workload (number of alerts, log scale) on the test period.}
\label{fig:tradeoff}
\end{figure}

\subsection{Cost-based operating point}
A simple cost model makes the trade-off concrete. Assume each alert costs a fixed review cost $c$, and each missed fraud
loses its transaction amount. Reviewing the top $k$ alerts then costs
\begin{equation}
\text{Cost}(k)=c\cdot k+\sum_{i\,\in\,\text{frauds not in top-}k}\text{Amount}_i .
\end{equation}
With $c=5$, the cost is minimised by reviewing the top \costBestAlerts{} alerts, for a total of \costTotal{}, compared
with \costNone{} if no model is used (Fig.~\ref{fig:cost}). The cost values are illustrative, but the method lets a
bank choose its operating point from its own review and loss costs.

\begin{figure}[H]
\centering
\includegraphics[width=0.8\linewidth]{02_cost_curve.png}
\caption{Review cost plus missed-fraud amount versus the number of alerts reviewed, for three review costs per alert.}
\label{fig:cost}
\end{figure}

\subsection{Stability: the high-dimensional gap, measured}\label{sec:stab}
Table~\ref{tab:stab} and Fig.~\ref{fig:stab} refit the GMM with five EM seeds. With standard scaling at $K=\selK$,
test PR-AUC ranges from \stabSmin{} to \stabSmax{} (std \stabSstd{}): the same model, started differently, is either
useful or useless. With quantile inputs the range narrows to \stabQmin{}--\stabQmax{} (std \stabQstd{}), and the
rankings and top-100 queues agree far more. Instability grows with $K$ for both input types. Even for the final
model, the fitted covariance matrices have condition numbers up to about \condMax{}, so several directions are almost
flat. This shows how close a \selD-dimensional full-covariance mixture is to degeneracy. Moderate regularisation
($\lambda$ between 0 and $10^{-2}$) barely changes validation PR-AUC; strong regularisation ($\lambda=0.1$) blurs the
density and hurts.

\begin{table}[H]
\centering\small
\caption{Seed stability on test (5 EM seeds): PR-AUC mean $\pm$ std and range, mean Spearman $\rho$ between seeds,
and mean overlap of the top-100 alerts.}
\label{tab:stab}
\begin{tabular}{@{}r l c c c c@{}}
\toprule
\textbf{$K$} & \textbf{Inputs} & \textbf{PR-AUC} & \textbf{Range} & \textbf{Spearman $\rho$} & \textbf{Top-100 overlap}\\
\midrule
\stabilityRows
\bottomrule
\end{tabular}
\end{table}

\begin{figure}[H]
\centering
\includegraphics[width=0.75\linewidth]{02_seed_stability.png}
\caption{Test PR-AUC of five EM seeds per configuration: quantile inputs (green) versus standard scaling (red).}
\label{fig:stab}
\end{figure}

\subsection{Error analysis}
At the 0.5\% budget the caught frauds have extreme values in V14 and V17 (mean $|\text{V14}|\approx8$), while the
missed frauds look ordinary in these components (about a quarter of that) and sit in dense legitimate regions. False
alarms are legitimate transactions in sparse regions: unusual component values and larger amounts. They are ``rare but
honest'' behaviour, which any purely density-based detector will flag.

%==================================================================
\section{Discussion}\label{sec:disc}

\subsection{Interpretation}
The results show that, for density-based fraud detection, the \emph{representation} matters more than the model's
flexibility. The quantile transform moves every feature onto a common Gaussian scale. It removes the heavy tails that
otherwise dominate the covariance estimates, and it makes the joint structure of legitimate behaviour close enough to
Gaussian that even one component works well. Information criteria answer a different question (``which model
describes legitimate data best?'') from the one we care about (``which model ranks fraud highest?''). At this sample
size they keep favouring larger mixtures that do not detect fraud better.

\subsection{The gap that motivates Phase 2}
\begin{enumerate}
\item \textbf{Shape assumption.} The GMM works only after every marginal is \emph{forced} to look Gaussian, and the
clipping in that transform discards information about how extreme a value is. The joint structure remains
non-Gaussian.
\item \textbf{Instability in high dimensions.} Each component costs \perComp{} parameters in \selD{} dimensions. EM's
solution depends on its seed, and the covariances are near-singular.
\item \textbf{Likelihood $\neq$ detection.} Unsupervised model selection (AIC/BIC) keeps asking for larger models that
detect fraud no better and are less stable.
\item \textbf{Mixture adds little.} $K=1$, $K=\selK$ and $K=\kBIC$ perform alike. Adding ellipsoids does not capture
the curved, low-dimensional structure of legitimate behaviour.
\end{enumerate}
Phase 2 replaces the fixed Gaussian shape with a \emph{learned non-linear} model of legitimate behaviour (undercomplete
and denoising autoencoders~\cite{sakurada2014,vincent2008}). It then tests hybrids that combine the Gaussian density
with the autoencoder, in the spirit of DAGMM~\cite{zong2018}.

\subsection{Limitations and threats to validity}
\begin{itemize}
\item \textbf{Short horizon.} The data covers two days in 2013; long-term concept drift is untested.
\item \textbf{Anonymised features.} PCA anonymisation prevents domain features (merchant, customer history, device).
\item \textbf{Daytime test period.} Validation and test fall in daytime hours (Section~\ref{sec:data}).
\item \textbf{Few test frauds.} With 74 test frauds, confidence intervals are about $\pm0.1$ PR-AUC wide; differences
smaller than this should not be over-interpreted.
\item \textbf{Development subsample.} Selection studies use a 40{,}000-row subsample; the final model uses all rows.
\end{itemize}

%==================================================================
\section{Conclusion and Future Work}\label{sec:concl}
A Gaussian-mixture density of legitimate transactions, fitted without any fraud label, ranks fraud well enough for
practical triage. It reaches a test PR-AUC of \gmmTestAP{}, catches \gmmRecallA{} of frauds at a 0.1\% false-positive
budget, and fills a top-50 review queue that is \gmmPA{} fraud. The decisive ingredient was feature engineering
(marginal Gaussianisation), not the number of components. The information criteria the textbook recommends keep asking
for larger models that bring no detection gain at this sample size. Together with the measured EM instability, these
findings define the gap addressed in Phase~2: learning a non-linear model of legitimate behaviour with autoencoders,
and testing whether a hybrid of the two paradigms can do better than either alone.

%==================================================================
\bibliographystyle{IEEEtran}
\bibliography{../references}

%==================================================================
\appendix
\section{Reproducibility}
\begin{itemize}
\item \textbf{Code.} Four self-contained notebooks (\texttt{notebooks/01}--\texttt{04}) run top-to-bottom on Google
Colab. Phase 1 uses notebooks 01 (data and features) and 02 (GMM). Each notebook downloads the dataset from a public
mirror and verifies its SHA-256 checksum.
\item \textbf{Seeds.} Global seed 42, plus explicit seed lists for multi-seed studies.
\item \textbf{Traceability.} Every number in this report is generated automatically from the notebooks' saved results
(\texttt{results/*.json}) into \texttt{numbers.tex}. Every figure is saved by the notebooks into \texttt{figures/}.
\item \textbf{Run time on Colab.} Notebook 01 takes about 2 minutes and notebook 02 about 35 minutes (CPU runtime).
\end{itemize}

\section{Team Contributions}
\begin{table}[H]
\centering\small
\caption{Team T068 (Section B): members and primary responsibilities. Every member contributed to the analysis and
can explain the full pipeline.}
\begin{tabularx}{\linewidth}{@{}l l X@{}}
\toprule
\textbf{Name} & \textbf{Enrollment No.} & \textbf{Primary responsibility (Phase 1)}\\
\midrule
Sharma Piyush & 2401010437 & Data quality, exploratory analysis, chronological split and leakage control\\
Saransh Singh Tomar & 2401010418 & Feature engineering and scaling ablation (quantile transform, cyclic time, log amount)\\
Satwik Mani Tripathi & 2401020056 & GMM theory, EM algorithm, AIC/BIC model selection and stability study\\
Gourav Tanwar & 2401010173 & Reference baselines (single Gaussian, Isolation Forest, supervised reference) and error analysis\\
Alok Kumar & 2401010058 & Evaluation protocol, threshold and cost analysis, report and presentation\\
\bottomrule
\end{tabularx}
\end{table}

\end{document}
